\documentclass[../../main.tex]{subfiles}
\begin{document}

{\bf [解]}

$f(x)$は3次式であり，条件(B)から，$a\neq0$として
\begin{align}
  f''(x)=2ax+2 \label{eq:1}
\end{align}
と書ける．積分して条件$f'(0)=0$を用いると
\begin{align}
  f'(x)=ax^2+2x \label{eq:2}
\end{align}
だから，さらに積分して実数$b$を用いて
\begin{align}
  f(x)=\dfrac{a}{3}x^3+x^2+b \label{eq:3}
\end{align}
と書ける．

一方で条件(A)より$f(x)-x$は$x=-t$，$x=0$を根にもつので
\begin{align}
  &
  \begin{cases}
    f(-t)+t = 0 \\
    f(0) = 0
  \end{cases} \\
  \therefore
  &
  \begin{cases}
    -\dfrac{a}{3}t^3+t^2+t+b=0 \\
    b=0
  \end{cases} \\
  \therefore
  &
  \begin{cases}
    t\left(-\dfrac{a}{3}t^2+t+1\right) = 0 \\
    b=0
  \end{cases} \label{eq:4}
\end{align}
である．$t>0$から第一式は
\begin{align}
  a=\dfrac{3(t+1)}{t^2} \label{eq:5}
\end{align}
となる．以上\cref{eq:4,eq:5}より$f(x)$は
\begin{align}
  f(x) = \dfrac{a}{3}x^3+x^2 \label{eq:6}
\end{align}
であり，$f(x)-x$を$g(x)$とおくと
\begin{align}
  g(x)
  &=f(x)-x \\
  &=\dfrac{t+1}{t^2}x^3+x^2-x\\
  &=\dfrac{1}{t^2}x(x+t)\Bigl((t+1)x-t\Bigr)
\end{align}
と因数分解できる．これを図示すると下図であり，求める面積$S(t)$は斜線部の領域の面積である．

\begin{figure}[htb]
  \centering
  \begin{tikzpicture}[scale=3.0, >=stealth, every node/.style={font=\small}]

    % --- 1. 定数・代表点のパラメータ定義 ---
    % t = 1 のとき:
    % a = 3*(1+1)/1^2 = 6  =>  f(x) = 2x^3 + x^2
    % d = t/(t+1) = 1/2 = 0.5
    % 交点: x = -1 (-t), x = 0 (原点), x = 0.5 (d)
    \def\tVal{1.0}
    \def\dVal{0.5}

    % --- 2. 求める面積 S(t) のハッチング領域 (0 <= x <= d) ---
    % 直線 y = x が上、曲線 y = f(x) が下
    \fill[pattern=north east lines, pattern color=blue!60, opacity=0.7]
    (0, 0) -- (\dVal, \dVal)
    -- plot[domain=\dVal:0, samples=60] (\x, {2.0*\x*\x*\x + \x*\x})
    -- cycle;

    % --- 3. 座標軸 ---
    \draw[->, thick] (-1.35, 0) -- (0.95, 0) node[right] {$x$};
    \draw[->, thick] (0, -1.35) -- (0, 1.15) node[above] {$y$};
    \node[below left=2pt] at (0, 0) {$O$};

    % --- 4. 直線 y = x ---
    \draw[thick, orange!90!black, domain=-1.25:0.85]
    plot (\x, {\x});
    \node[orange!90!black, font=\footnotesize, above right] at (0.75, 0.75) {$y = x$};

    % --- 5. 3次曲線 y = f(x) ---
    \draw[very thick, blue!85!black, domain=-1.18:0.65, samples=100]
    plot (\x, {2.0*\x*\x*\x + \x*\x});
    \node[blue!85!black, font=\footnotesize, left] at (0.62, {2.0*0.62^3 + 0.62^2}) {$y = f(x)$};

    % --- 6. 交点および射影補助線 ---
    % (i) 交点 P(-t, -t)
    \draw[dashed, gray!70] (-\tVal, 0) -- (-\tVal, -\tVal) -- (0, -\tVal);
    \fill[red!80!black] (-\tVal, -\tVal) circle (0.7pt);
    \node[above left=1pt, font=\footnotesize, red!80!black] at (-\tVal, -\tVal) {$P$};
    \node[above=2pt, font=\scriptsize] at (-\tVal, 0) {$-t$};
    \fill (-\tVal, 0) circle (0.5pt);

    % (ii) 原点 O(0, 0)
    \fill[black] (0, 0) circle (0.6pt);

    % (iii) 交点 Q(d, d)
    \draw[dashed, gray!70] (\dVal, 0) -- (\dVal, \dVal) -- (0, \dVal);
    \fill[red!80!black] (\dVal, \dVal) circle (0.7pt);
    \node[above left=1pt, font=\footnotesize, red!80!black] at (\dVal, \dVal) {$Q$};
    \node[below=2pt, font=\scriptsize] at (\dVal, 0) {$d = \frac{t}{t+1}$};
    \fill (\dVal, 0) circle (0.5pt);

  \end{tikzpicture}
  \caption{3次曲線 $y = f(x)$ と直線 $y = x$，および面積 $S(t)$ を与える領域（斜線部）}
  \label{fig:cubic_line_area_St}
\end{figure}

簡単のために$d=t/(t+1)$とおくと求める面積$S(t)$は
\begin{align}
  S(t)
  &=-\int_0^dg(x)\,dx\\
  &=\Bigl[-\dfrac{t+1}{4t^2}x^4-\dfrac{1}{3}x^3+\dfrac{1}{2}x^2\Bigr]_0^d\\
  &=-\dfrac{1}{4}\Bigl(\dfrac{1}{t}+\dfrac{1}{t^2}\Bigr)
  \Bigl(\dfrac{1}{1+1/t}\Bigr)^4
  -\dfrac{1}{3}\Bigl(\dfrac{1}{1+1/t}\Bigr)^3
  +\dfrac{1}{2}\Bigl(\dfrac{1}{1+1/t}\Bigr)^2
\end{align}
となる．よって$t\to\infty$の極限をとって
\begin{align}
  \lim_{t\to\infty}S(t)=\dfrac{1}{2}-\dfrac{1}{3}=\dfrac{1}{6}
\end{align}
である．

\newpage
\end{document}
